% 原书第 8.7 节。
\section{推理}
\label{sec:ch08-inference}
为大语言模型应用设计服务系统时，需要就如何执行模型推理作出多项决定。大部分决定涉及成本、速度和模型表现三个方面的优化。700 亿参数的模型可以提供质量很高的输出，但其速度远低于 70 亿参数的模型，除非愿意在计算基础设施上投入大量资金。本节介绍在这些方面进行权衡的方法，以及因使用场景而异的若干重要因素。

\subsection{模型托管}
大语言模型应用的一项基本决策是模型本身托管在哪里。可以把选项分为以下三类。

\begin{itemize}
\item \textbf{向公开的第三方 API 发送推理请求。}这种方式快捷、直接，常用于构建演示和原型，以及迅速将新构想投入生产。它无需自行搭建或维护，开发人员无需深入了解大语言模型的原理就可学会使用这些 API。部署应用所需的工作和专业知识较少，可以节省不少费用。不过，API 本身的使用成本相对较高，推理请求量大时，很容易抵消所节省的开销。还有几个限制需要考虑。首先，这种方式几乎无法针对特定使用场景微调或以其他方式定制模型。其次，延迟没有强保证，还须考虑公开 API 常见的速率限制。最后，如第~6.4~节所述，这些 API 调用意味着应用中的数据会交给第三方。对许多组织来说，仅这一点就足以放弃该方案。尽管有多种缺点，OpenAI 当时最新的 GPT 模型只通过其 API 提供。如果采用 OpenAI 当时领先的模型能力可以为终端用户提供最大的价值，这些取舍可能是值得的。

\item \textbf{使用基础模型托管服务。}三家主要云计算提供商均有相应服务，例如 AWS Bedrock，可在私有、安全的环境中使用基础模型。出于多种原因，这种方式比公开 API 更容易扩展。首先，服务提供商为推理提供多项内置优化，账号所有者也可以在一定程度上控制分配给模型的 GPU 资源，从而在推理速度和计算开销之间找到合适的平衡。API 用户无法控制这两项因素。其次，如果大语言模型与应用的其他部分位于同一个云环境，许多网络问题可以减轻。推理请求和响应较少受带宽波动影响，也可以通过配置让模型与应用位于同一物理地点。延迟很重要时，把请求发送到可能远在地球另一端的 API 会带来问题。不过，模型托管服务也有限制。为了优化算力使用，服务所支持的基础模型和调优技术往往是预先规定的一组。优化还附带额外费用，对于内部已有 GPU 计算专长的组织，这笔开销未必划算。

\item \textbf{在自己的计算基础设施上托管模型。}如果 AWS Bedrock 等服务限制过多，可以自行搭建运行环境。这样可以灵活选择大语言模型，并针对应用需求精确优化，复杂度也高于前两种方案。NVIDIA 的 Triton 推理服务器\footnote{\url{https://github.com/triton-inference-server}}可以减少一部分工作。它支持灵活选择模型架构，并负责多项底层 GPU 优化。许多组织无需聘用或外包一支熟悉张量运算的专门团队，也能构建高度定制的推理系统。不过，推理需求达到很大规模时，这种投入可能带来成本优势。
\end{itemize}

\subsection{性能优化}
\label{sec:ch08-inference-performance}
在 CPU 上尝试过运行大语言模型的人，多半体会过 GPU 算力不足时的响应速度。算力成本较高，人们因此提出了多种方法，以便在不增加硬件的情况下更快地处理更多推理请求。

与推理速度有关的两个重要概念是延迟和吞吐量。延迟指处理一个请求并向应用发送响应所需的时间，吞吐量指大语言模型在给定时间内可以产生的输出量（\refcite{ch08-ref-001}{Agarwal et al., 2023}）。二者关系密切，但含义不同。例如，用户开始输入代码时，期待编程助手几乎立即给出建议。这类应用要求低延迟。再设想一个过滤垃圾邮件的服务。无论分类花半秒、几秒还是更久，用户很可能感觉不到区别，但吞吐量仍然重要。如果服务跟不上新邮件的涌入，积压就会越来越多，无法发挥预期作用。

\subsubsection{批处理}
\label{sec:ch08-batching}
优化批处理对于提高吞吐量很重要，见图~\ref{fig:ch08-batching}。把多个输入合入同一个矩阵，可以利用向量化带来的性能收益。一般来说，较大的批次比小批次效率更高。如果应用的大多数推理请求包含大量输入数据，批处理容易组织。若用户每次只发送一个输入，例如在对话应用中，则要把多个用户的输入合入同一批次，才能利用批处理。这称为动态批处理（dynamic batching）。它能提高吞吐量，但用户可能需要等待服务器积攒其他请求后才开始处理自己的请求，从而增大延迟。

连续批处理（continuous batching）是针对自回归大语言模型解决上述问题的方法。生成式文本模型迭代地产生词元，因而可以把新的用户输入加入正在处理的批次。若某个输出词元串已完成，即达到了长度上限或生成了停止词元，就会腾出批次中的一个输入位置。随后可以把队列中的下一个用户请求放入批次。这样，GPU 内存一旦腾出空间，系统就能开始处理新的请求，也能避免因批次小于合适规模而未充分利用 GPU。这种方法也能处理长度相差很大的输入，而不必为额外的填充词元付出代价。GPU 专长于大型矩阵运算，输入的大小和形状持续适合硬件架构时，其性能最高。

\sourcefigure{fig08_06.pdf}{标准批处理、动态批处理和连续批处理。传统推理架构主要由客户端组成批次。尤其在每个用户每次只发一个请求的应用中，服务器动态汇聚多个输入能更充分地利用 GPU，但凑齐批次后才能开始计算，因而会增加等待时间。连续批处理把新收到的输入加入正在执行的批次，以减少这种等待。}{fig:ch08-batching}

\subsubsection{键值缓存}
键值缓存（key-value caching）是另一种适用于自回归大语言模型的推理技术。每生成一个词元，就把它接在序列末尾，再将序列送入模型，以生成下一个词元。此前所有词元都与上一轮相同，因此其对应的键和值张量可缓存并反复使用，无需在每轮迭代中重新计算。这样只需为新词元生成键和值。\footnote{原文把可复用的量称为注意力权重。键值缓存存储的是此前词元的键和值，新词元仍需计算相对于这些键的注意力权重。}注意力机制是 Transformer 中时间复杂度最高的部分，通常也是架构中最大的性能瓶颈。减少这些计算可以显著提高推理速度。

\subsection{成本优化}
即便采用所有可用的推理加速技术，规模最大、能力最强的模型仍然需要价格较高的硬件。应用同时要求低延迟和高吞吐量时尤其如此。几乎总是值得考察较小的模型能否同样好地完成工作，或至少取得相近的表现。在某些使用场景中，限制客户选择的可能更多是预算，而非结果质量。降低成本与提高效用之间存在一系列可行取舍，生产应用都需要仔细分析，在其中找出合适的平衡点。

模型规模与结果质量之间还有一个可权衡的方面。采用较小模型节省的一部分费用可用来微调，以弥补能力差距。大模型的吸引力之一在于它包含丰富的知识，仅用提示工程和上下文学习，就能较好地完成广泛的任务。这很重要，因为即使采用 LoRA 等技术，微调大模型也很昂贵。选择较小的大语言模型后，微调更可行。Predibase 的业务建立在这一想法上。其主张是，多数应用适合使用较小的专用模型，并按各类推理请求的需要配备适配器。
